\documentclass[aspectratio=169, 12pt]{beamer}

\usetheme{Madrid}
\usecolortheme{whale}

\definecolor{azulprincipal}{RGB}{10, 40, 90}
\definecolor{azulsecundario}{RGB}{30, 90, 170}
\definecolor{acento}{RGB}{200, 50, 50}
\definecolor{grisclaro}{RGB}{240, 240, 245}

\setbeamercolor{palette primary}{bg=azulprincipal, fg=white}
\setbeamercolor{palette secondary}{bg=azulsecundario, fg=white}
\setbeamercolor{palette tertiary}{bg=azulprincipal, fg=white}
\setbeamercolor{palette quaternary}{bg=azulprincipal, fg=white}
\setbeamercolor{structure}{fg=azulprincipal}
\setbeamercolor{block title}{bg=azulprincipal, fg=white}
\setbeamercolor{block body}{bg=grisclaro, fg=black}
\setbeamercolor{block title alerted}{bg=acento, fg=white}
\setbeamercolor{block body alerted}{bg=acento!10, fg=black}
\setbeamercolor{frametitle}{bg=azulprincipal, fg=white}


% Paquetes

\usepackage[utf8]{inputenc}
\usepackage[T1]{fontenc}
\usepackage{amsmath, amssymb}
\usepackage{mathtools}
\usepackage{physics}
\usepackage{booktabs}
\usepackage{array}
\usepackage{xcolor}
\usepackage{listings}
\usepackage{tikz}
\usetikzlibrary{arrows.meta, positioning, shapes.geometric}


% Configuración de listings para Julia

\lstdefinelanguage{Julia}{
  keywords={function, end, for, if, else, return, using, const, in, do, while},
  keywordstyle=\color{azulprincipal}\bfseries,
  comment=[l]{\#},
  commentstyle=\color{gray}\itshape,
  stringstyle=\color{acento},
  morestring=[b]",
  sensitive=true
}

\lstset{
  language=Julia,
  basicstyle=\ttfamily\footnotesize,
  backgroundcolor=\color{grisclaro},
  frame=single,
  framerule=0pt,
  rulecolor=\color{azulprincipal},
  breaklines=true,
  showstringspaces=false,
  tabsize=2
}

% Comandos propios

\newcommand{\that}[1]{\hat{#1}}
\newcommand{\F}{\mathcal{F}}
\newcommand{\uv}{\mathbf{u}}


% Información de la presentación

\title[Kelvin-Helmholtz Pseudo-Espectral]{%
  Simulación Pseudo-Espectral de la\\
  Inestabilidad de Kelvin-Helmholtz
}
\author{Pedro Damian Meneses Orozco}
\institute{DMD}
\date{\today}



\begin{document}


\begin{frame}
  \titlepage
\end{frame}
% Índice

\begin{frame}{Contenido}
  \tableofcontents
\end{frame}


\section{Inestabilidad de Kelvin-Helmholtz (IKH)}


\begin{frame}{¿Qué es una inestabilidad de Kelvin-Helmholtz?}
  \begin{columns}
    \begin{column}{0.55\textwidth}
      Este fenómeno es una inestabilidad hidrodinámica que aparece cuando dos fluidos adyacentes tienen una velocidad relativa en su interfaz, o en otro caso, cuando hay cizallamiento de velocidad de un único fluido, este segundo caso es el que estudiaremos en esta presentación.

      El motivo del interés por este tema es que sigue siendo muy estudiado en general por su presencia en la naturaleza.
      Por ejemplo:
.
    \end{column}
    \begin{column}{0.42\textwidth}
      \begin{figure}[H]
        \centering
        \includegraphics[width=0.8\textwidth]{Saturn_Kelvin_Helmholtz.jpeg}
        \caption{Una IKH en Saturno, por la interacción de dos bandas en la atmósfera del planeta (foto de la sonda Cassini.  NASA/JPL/Space Science Institute).}
        \label{fig:grafica}
        \end{figure}
    \end{column}
  \end{columns}
\end{frame}

\begin{frame}
  En si, este tipo de inestabilidades ofrecen una posible transición entre un flujo laminar y uno turbulento, al a ver efectos de la viscosidad en la interfaz. veremos que ocurre con nuestro caso. Para nuestro caso simplificaremos las ecuaciones a resolver: tomaremos un fluido bidimensional, newtonieano e incompresible. Entonces el campo de velocidades debe satisfacer:

  \vspace{0.3em}
  \begin{columns}
    \begin{column}{0.48\textwidth}
      \begin{block}{Navier-Stokes incompresible 2D}
        \[
          \pdv{\vec{u}}{t} + (\vec{u}\cdot\nabla)\vec{u} =
          -\frac{\nabla p}{\rho} + \nu\nabla^2\vec{u}
        \]
        \[
          \nabla\cdot\uv = 0
        \]
      \end{block}

      \vspace{0.5em}
      Con $\nabla = (\frac{\partial}{\partial x},\frac{\partial}{\partial y} )$. Notemos que queremos hacer una simulación numérica y queremos ver lo que ocurre en el cizallamiento por lo que lo analizaremos la vorticidad $\omega = \nabla \times \vec{u} \cdot \hat{e_z}$
    \end{column}
    \begin{column}{0.48\textwidth}
       pues al ser un fluido en 2 dimensiones unicamente existirá componente en z. Sustituyendo eso en la ecuación de Navier-Stokes y usando la ecuación de incompresibilidad llegamos a:
       \begin{alertblock}{Ecuación de vorticidad}
        \[
          \pdv{\omega}{t} + u\pdv{\omega}{x} + v\pdv{\omega}{y}
          = \nu\nabla^2\omega
        \]
      \end{alertblock}
    \end{column}
  \end{columns}
\end{frame}

\begin{frame}
  \begin{columns}
    \begin{column}{0.48\textwidth}
      Nos deshicimos del gradiente de presión, lo que nos ahorrará computo.

      Ahora, por las clases sabemos que para cualquier campo incompresible en dos dimensiones $\exists\psi $ función de corriente, tal que, si $\vec{u} = (u, v)$:

    $$u =  \frac{\partial\psi}{\partial y} $$
    $$v = - \frac{\partial\psi}{\partial x} $$
    Y entonces la vorticidad:
    $$\omega = \frac{\partial v}{\partial x} - \frac{\partial u}{\partial y} = -\frac{\partial^2 \psi}{\partial x^2} - \frac{\partial^2 \psi}{\partial y^2} = - \nabla^2\psi$$
    \end{column}
    \begin{column}{0.48\textwidth}
      \begin{block}{Esto convierte...}
        A la vorticidad en la solución  de una ecuación de Poisson:
        $$\nabla^2\psi=-\omega$$
      \end{block}
      Con esto, ya podemos solucionar todo:
      De $\omega$ obtenemos $\psi$ de la ecuación de Poisson, luego de $\psi$ obtenemos (u,v), esto nos regresa a la ecuación de vorticidad e integrando respecto al tiempo obtenemos $\omega$. 
    \end{column}
  \end{columns}
\end{frame}

\section{Formulación Matemática}

\begin{frame}{¿Qué es un método espectral?}
  \begin{columns}
    \begin{column}{0.55\textwidth}
      En lugar de discretizar el dominio en celdas pequeñas,
      la solución se expresa como suma de funciones base globales:
 
      \vspace{0.5em}
      \begin{block}{Representación espectral}
        \[
          u(x,t) \approx \sum_k \hat{u}_k(t)\,\phi_k(x)
        \]
      \end{block}
 
      \vspace{0.5em}
      donde $\phi_k$ son los \textbf{modos} y $\hat{u}_k$ los
      \textbf{coeficientes espectrales}.
    \end{column}
    \begin{column}{0.42\textwidth}
      \begin{block}{Ventaja clave}
        \textbf{Convergencia exponencial}: el error cae como
        $e^{-N}$ al aumentar el número de modos $N$, lo que ofrece alta precisión.
      \end{block}
 
      \vspace{0.5em}
      \begin{alertblock}{Limitación}
        Requieren geometrías simples (periódicas o rectangulares).
      \end{alertblock}
    \end{column}
  \end{columns}
\end{frame}

\begin{frame}{Transformada de Fourier Discreta (DFT)}
Usaremos para solucionar nuestras ecuaciones la DFT:

Sea $f(x)$ una función periódica de período $L$ discretizada en
$N$ puntos $x_j = jL/N$, $j = 0,\ldots,N-1$.
La DFT la descompone como:
 

  $$f(x_j) = \sum_{k=-N/2}^{N/2-1} \that{f}_k\,e^{ikx_j}$$

 
donde los coeficientes espectrales son:

  $$\that{f}_k = \frac{1}{N}\sum_{j=0}^{N-1} f(x_j)\,e^{-ikx_j}$$

El índice $k$ es el \textbf{número de onda}: indica cuántos ciclos
completos de la onda $e^{ikx}$ caben en el dominio $[0,L)$.
El máximo número de onda representable es $k_{\max} = N/2$
(\textbf{frecuencia de Nyquist}).

  
\end{frame}

\begin{frame}{Propiedad fundamental que usaremos de la transformada de Fourier}
La propiedad más poderosa de la representación de Fourier es que
\textbf{derivar en espacio físico equivale a multiplicar en espacio
espectral}:
 
\begin{equation}
  \widehat{(\pdv{f}{x}})_k = ik\,\that{f}_k.
  \label{eq:deriv_espectral}
\end{equation}
 
En dos dimensiones, el Laplaciano se convierte en:
 
 $$ \widehat{\nabla^2 f}_{\vec{K}} = -(k_x^2 + k_y^2)\that{f}_{\vec{K}} = -K^2\,\that{f}_\vec{K}$$
 
donde $K^2 = k_x^2 + k_y^2$ es el cuadrado del módulo del
vector de onda $\vec{K} = (k_x, k_y)$.

\end{frame}

\begin{frame}{La ecuación completa en espacio espectral}
  Aplicando la DFT a la ecuación de vorticidad:

  \begin{alertblock}{Ecuación espectral central}
    \[
      \pdv{\hat{\omega}_\mathbf{k}}{t} =
      \underbrace{-\widehat{(u\,\partial_x\omega
                  + v\,\partial_y\omega)}_\mathbf{k}}_{\text{advección (termino no lineal)}}
      \underbrace{- \nu K^2\,\hat{\omega}_\mathbf{k}}_{\text{difusión (exacta)}}
    \]
  \end{alertblock}

  \vspace{0.5em}
  \begin{columns}
    \begin{column}{0.48\textwidth}
      \begin{block}{Resolver Poisson en espectral}
        \[
          \hat{\psi}_\mathbf{k} = \frac{\hat{\omega}_\mathbf{k}}{K^2}
        \]
        Una división. Sin sistema lineal.
      \end{block}
    \end{column}
    \begin{column}{0.48\textwidth}
      \begin{block}{Velocidades en espectral}
        \[
          \hat{u_k} = ik_y\hat{\psi_k}, \quad
          \hat{v_k} = -ik_x\hat{\psi_k}
        \]
        Una multiplicación por número complejo.
      \end{block}
    \end{column}
  \end{columns}
\end{frame}

% ============================================================
\section{Método Pseudo-Espectral}
% ============================================================

\begin{frame}{El problema del término no lineal}
  El término advectivo $u\,\partial_x\omega + v\,\partial_y\omega$
  es un \textbf{producto} de dos campos. En espectral, un producto
  es una \textbf{convolución}:
  \[
    \widehat{fg}_k = \sum_{p+q=k} \hat{f}_p\,\hat{g}_q
    \quad \Rightarrow \quad \mathcal{O}(N^4) \text{ en 2D}
  \]

  \vspace{0.3em}
  \begin{alertblock}{Estrategia pseudo-espectral (Orszag, 1972)}
    \begin{enumerate}
      \item Calcular derivadas en espectral.
      \item Transformar al físico con IFFT.
      \item Calcular el producto en físico.
      \item Regresar al espectral con FFT.
    \end{enumerate}
    \vspace{0.2em}
    \textbf{Coste total:} $\mathcal{O}(N^2\log N)$ en lugar de $\mathcal{O}(N^4)$
  \end{alertblock}
\end{frame}

\begin{frame}{El aliasing}
  \begin{columns}
    \begin{column}{0.55\textwidth}
      Con $N$ puntos el número de onda máximo es $k_{\max} = N/2$.

      \vspace{0.5em}
      El producto de dos campos con modos hasta $k_{\max}$ genera
      modos hasta $2k_{\max} = N$. Estos modos no son
      representables en la malla y se ``doblan'' de vuelta:

      \[
        k = k_{\max} + m \;\longleftrightarrow\; k = k_{\max} - m
      \]

      El modo $k=5$ con $N=8$ es indistinguible del modo $k=3$
      en todos los puntos de la malla.
    \end{column}
    \begin{column}{0.42\textwidth}
      \begin{block}{Regla de los 2/3 (Orszag, 1971)}
        Truncar los modos a $\frac{2}{3}k_{\max}$ antes de
        cada producto:
        \[
          \hat{f}_\mathbf{k}^{\text{trunc}} =
          \begin{cases}
            \hat{f}_\mathbf{k} & |\mathbf{k}| \leq \tfrac{2}{3}k_{\max}\\
            0 & \text{si no}
          \end{cases}
        \]
        Entonces los productos generan modos hasta
        $\frac{4}{3}k_{\max} < 2k_{\max}$. Sin aliasing.
      \end{block}
    \end{column}
  \end{columns}
\end{frame}


% ============================================================
\section{Integración Temporal: RK4}
% ============================================================

\begin{frame}{Runge-Kutta de orden 4}
  La ecuación tiene la forma abstracta $\frac{d\omega}{d t} = \mathcal{F}(\omega)$. Donde $\mathcal{F}$ es todo el procedimiento que hemos hecho para encontrar la derivada temporal.
  RK4 evalúa $\mathcal{F}$ cuatro veces por paso:

  \begin{columns}
    \begin{column}{0.5\textwidth}
      \begin{block}{Las cuatro etapas}
        \begin{align*}
          k_1 &= \mathcal{F}(\omega^n)\\[4pt]
          k_2 &= \mathcal{F}\!\left(\omega^n + \tfrac{\Delta t}{2}k_1\right)\\[4pt]
          k_3 &= \mathcal{F}\!\left(\omega^n + \tfrac{\Delta t}{2}k_2\right)\\[4pt]
          k_4 &= \mathcal{F}(\omega^n + \Delta t\,k_3)
        \end{align*}
      \end{block}
    \end{column}
    \begin{column}{0.46\textwidth}
      \begin{alertblock}{Actualización}
        \[
          \omega^{n+1} = \omega^n +
          \frac{\Delta t}{6}(k_1 + 2k_2 + 2k_3 + k_4)
        \]
      \end{alertblock}

    \end{column}
  \end{columns}

\end{frame}

% ============================================================
\section{Inestabilidad de Kelvin-Helmholtz}
% ============================================================

\begin{frame}{El perfil de cizallamiento}
  \begin{columns}
    \begin{column}{0.55\textwidth}
      Dos capas de fluido a velocidades opuestas, separadas por
      una capa de mezcla de espesor $\delta$:

      \begin{block}{Perfil tangente hiperbólica}
        \[
          U(y) = U_0\tanh\!\left(\frac{y}{\delta}\right)
        \]
      \end{block}

      \vspace{0.4em}
      \textbf{¿Por qué tanh?}
      \begin{itemize}
        \item Es la solución física exacta por difusión viscosa
        \item Infinitamente diferenciable: preserva la convergencia espectral
        \item $\delta$ controla el cizallamiento: $dU/dy\big|_{y=0} = U_0/\delta$
      \end{itemize}
    \end{column}
    \begin{column}{0.42\textwidth}
      \begin{block}{Vorticidad base}
        \[
          \omega_0(y) = -\dv{U}{y}
          = -\frac{U_0}{\delta}\operatorname{sech}^2\!\!\left(\frac{y}{\delta}\right)
        \]
      \end{block}

      \vspace{0.3em}
      Campana centrada en $y=0$ — máximo cizallamiento en la interfaz.

      \vspace{0.4em}
      \begin{alertblock}{Condición inicial}
        \[
          \omega' = \varepsilon\cos(k^*x)\operatorname{sech}^2\!\!\left(\frac{y}{\delta}\right)
        \]
        $\varepsilon = 0.01 \ll 1$: régimen lineal inicial.
      \end{alertblock}
    \end{column}
  \end{columns}
\end{frame}

\begin{frame}{Análisis de estabilidad lineal}
  Perturbando el estado base $\omega = \omega_0 + \omega'$ y
  buscando modos normales $\omega' \sim \hat{\omega}'(y)e^{i(kx-\sigma t)}$:

  \begin{block}{Ecuación de Rayleigh}
    \[
      \left(U - \frac{\sigma}{k}\right)
      \left(\frac{d^2\hat{\omega}'}{dy^2} - k^2\hat{\omega}'\right)
      - \frac{d^2U}{dy^2}\,\hat{\omega}' = 0
    \]
  \end{block}

  \vspace{0.4em}
  \begin{columns}
    \begin{column}{0.48\textwidth}
      \begin{block}{Criterio de Rayleigh (1880)}
        Condición \emph{necesaria} para inestabilidad:
        el perfil $U(y)$ debe tener un punto de inflexión
        ($U'' = 0$ en algún $y_c$).

        \vspace{0.3em}
        El perfil tanh satisface esto en $y=0$.
      \end{block}
    \end{column}
    \begin{column}{0.48\textwidth}
      \begin{alertblock}{Resultado para el perfil tanh}
        \begin{itemize}
          \item Inestabilidad para todo $k < 1/\delta$
          \item Modo más inestable: $k^* \approx 0.45/\delta$
          \item Tasa de crecimiento $\sim \text{Im}(\sigma) > 0$
        \end{itemize}
      \end{alertblock}
    \end{column}
  \end{columns}
\end{frame}

% ============================================================
\section{Diagnósticos}
% ============================================================

\begin{frame}{Energía cinética y enstrofía}
  \begin{columns}
    \begin{column}{0.48\textwidth}
      \begin{block}{Energía cinética}
        \[
          E(t) = \frac{1}{2}\int_\Omega |\uv|^2\,dA
               = \frac{1}{2}\sum_\mathbf{k} K^2|\hat{\psi}_\mathbf{k}|^2
        \]
        Mide la intensidad total del flujo.
        Conservada si $\nu = 0$.
      \end{block}
    \end{column}
    \begin{column}{0.48\textwidth}
      \begin{block}{Enstrofía}
        \[
          Z(t) = \frac{1}{2}\int_\Omega \omega^2\,dA
               = \frac{1}{2}\sum_\mathbf{k} |\hat{\omega}_\mathbf{k}|^2
        \]
        Mide la intensidad total de rotación.
        Cuasi-conservada en turbulencia 2D.
      \end{block}
    \end{column}
  \end{columns}

  \vspace{0.5em}

\end{frame}



\begin{frame}{Referencias}
  \begin{enumerate}
    \item Orszag, S.A. (1971). \textit{On the elimination of aliasing
          in finite-difference schemes by filtering high-wavenumber
          components.} J. Atmos. Sci..

    \item Orszag, S.A. (1972). \textit{Comparison of pseudospectral
          and spectral approximations.} Stud. Appl. Math.
        

    \item Kraichnan, R.H. (1967). \textit{Inertial ranges in
          two-dimensional turbulence.} 

    \item Canuto et al. (1988). \textit{Spectral Methods in Fluid Dynamics.} Springer.

    \item Rayleigh, Lord (1880). \textit{On the stability, or
          instability, of certain fluid motions.}
          Proc. London Math. 

    \item Drazin \& Reid (2004). \textit{Hydrodynamic Stability}
          (2nd ed.). Cambridge University Press.
  \end{enumerate}
\end{frame}

% ============================================================
% Diapositiva final
% ============================================================
\begin{frame}
  \begin{center}
    \vspace{1cm}
    {\Large\bfseries\color{azulprincipal}
      Simulación Pseudo-Espectral de\\
      Kelvin-Helmholtz en Julia
    }

    \vspace{1cm}
    \textcolor{gray}{\rule{0.5\textwidth}{0.4pt}}
    \vspace{0.5cm}

    {\normalsize
      Métodos Espectrales $\cdot$
      Turbulencia 2D $\cdot$
      Julia / FFTW
    }
  \end{center}
\end{frame}

\end{document}